\documentclass[12pt]{article}

\usepackage[onehalfspacing]{referee-response}

\begin{document}

% Title ------------------------------------------------------------------------
\begin{center}
    {\Large\bf Paper Title}
    
    \vspace{1em}

    {\large Response to the Editor and Reviewers}

    \vspace{.5em}

    \emph{Journal} (\texttt{Review Number})
    
    \vspace{1em}
\end{center}
% ------------------------------------------------------------------------------

\noindent We thank the editor and the reviewers for the constructive and thoughtful comments. We have carefully considered each comment, and we have revised the manuscript accordingly. 
We believe these revisions have significantly improved the manuscript. In this response memorandum, we elaborate in detail on how we addressed each of the comments and questions of the editor and reviewers. For reference, we provide a brief summary of some of the main changes to the manuscript below.

\section{Summary of Revisions}

\begin{description}
  \item[Point 1] \lipsum[1]
  \item[Point 2] \lipsum[1]
\end{description}

% ------------------------------------------------------------------------------
\vspace*{\bigskipamount}\NewRef{Editor}{Editor}
% ------------------------------------------------------------------------------


\begin{refcomment}
  Full text of the referee's comment. It will automatically create a \texttt{\textbackslash subsection} with the comment number.
  
  \lipsum[1-2]
\end{refcomment}

Response.
\lipsum[2]



% ------------------------------------------------------------------------------
\vspace*{\bigskipamount}\NewRef{Referee 1}{R1}
% ------------------------------------------------------------------------------

\begin{refcomment}
  Full text of the referee's comment. It will automatically create a \texttt{\textbackslash subsection} with the comment number.
  
  \lipsum[1-2]
\end{refcomment}

Response.
\lipsum[3-4]

\begin{refcomment}
  Comment 
\end{refcomment}

Response

\begin{refcomment*}
  Comment 
\end{refcomment*}

% ------------------------------------------------------------------------------
\vspace*{\bigskipamount}\NewRef{Referee 2}{R2}
% ------------------------------------------------------------------------------

\begin{refcomment}[Comment 1a]
  You can use manual labels as an optional argument by 
  \begin{verbatim}
    \begin{refcomment}[Comment 1a] ... \end{refcomment}
  \end{verbatim}
\end{refcomment}
Response

\begin{refcomment}[Comment 1b]
  Comment 
\end{refcomment}
Response

\begin{refcomment}
  Comment 
\end{refcomment}
Response


% ------------------------------------------------------------------------------
\newpage
\section*{How to use this template}
% ------------------------------------------------------------------------------

Use \texttt{\textbackslash NewRef\{Name\}\{Abbreviation\}} to start each editor or referee section. The \texttt{Name} appears in the section heading, and the \texttt{Abbreviation} labels numbered comments.

Write comments with \texttt{\textbackslash begin\{refcomment\} ... \textbackslash end\{refcomment\}}. Comments are numbered automatically; use an optional argument, such as \texttt{\textbackslash begin\{refcomment\}[Comment 1a]}, to set a manual label. Use \texttt{refcomment*} for an unnumbered comment box. Responses should follow the comment block.


\end{document}
